\section{The Good, the Catastrophic, and the Extinguished Cycle}
\label{p11:sec:regimes}

\noindent
The paper has traversed the moments of a single turn. It now considers the cycle as a whole, in the regimes its turning may assume across the history of a relation. Three regimes are distinguished, by the conjunction of two questions: whether the cycle persists, and, if it persists, what sign of phase it accumulates. The two questions are not of the same order. The first concerns the existence of the relation, the second its character; and the regimes they jointly define are the flourishing spiral, the catastrophic cycle, and extinction.

\subsection{Persistence and phase as two questions}
\label{p11:sec:reg-two}

The dynamical vocabulary of \xref{p11:sec:dynsys} distinguishes two features of a cyclic system that bear on a relation. The first is whether the system sustains its cyclic movement at all, or decays toward the terminal rest of a fixed point. The second, for a system that does sustain its movement, is the phase it accumulates around a turn, positive, vanishing, or negative. These are independent: a cycle may persist with a positive phase, persist with a negative phase, or fail to persist; and the three regimes of a relation are the three outcomes of this pair of questions, the fourth combination, a non-persisting cycle with a phase, being empty, since a cycle that does not persist accumulates no phase.

\begin{claim}[Three regimes from two questions]
The regime of a relation is fixed by two independent questions: whether its cycle persists, and, if it persists, what sign of phase it accumulates over a turn. A persisting cycle of positive phase is a flourishing spiral, which ascends as it turns. A persisting cycle of negative phase is a catastrophic cycle, which persists while depleting, descending as it turns. A cycle that does not persist is extinction, the decay of the relation's movement toward the terminal rest in which value is no longer generated or returned. The question of persistence is prior to the question of phase, since a cycle that does not persist has no phase; the gravest condition of a relation is therefore not the negative phase but the failure of persistence, and the catastrophic cycle, for all its depletion, is a living relation in a way that extinction is not.
\end{claim}

\subsection{The flourishing spiral}
\label{p11:sec:reg-flourishing}

The flourishing spiral is the regime of the good cycle. Its value is generated and returned, its recyclability sustains the cycle above the threshold of decay, its return distribution is, relative to the coupling, approximately symmetric, and the unfolding of each party is sustained rather than foreclosed. Across its turns it accumulates a positive phase: the parties' cognition of the good deepens, their practice enriches, and they develop, through the shared life, into more than they were. The flourishing spiral is the regime the normative specification of \xref{p11:sec:cog-spec} describes, and it is, by the argument to follow, necessarily a just regime: the symmetry of its return distribution is the condition of its positive phase, and a relation whose return is severely skewed cannot sustain the ascent of all its parties, since the party consigned to the tail is depleted rather than developed by the turning.

\subsection{The catastrophic cycle}
\label{p11:sec:reg-catastrophic}

The catastrophic cycle is the regime of the bad cycle that nonetheless persists. Its movement is sustained, its form is intact, and it may score well on aggregate satisfaction; but it accumulates a negative phase, depleting its parties, or one of them, turn upon turn. This is the regime in which the orthogonal instruments of \xref{p11:sec:evaluation} do their characteristic work, for it is the regime a satisfaction measure cannot distinguish from the flourishing spiral, both persisting and both capable of a high mean, while the skew of return and the concentration of foreclosed undertakings distinguish them, the catastrophic cycle exhibiting the skew and the foreclosure that the flourishing spiral does not. The catastrophic cycle is the marriage that endures and depletes, the form maintained while one party, or both, is consumed in the maintaining; and it is the regime in which the injustice of a relation is most consequential, since its persistence prolongs the extraction that its intact form conceals.

\subsection{Extinction}
\label{p11:sec:reg-extinction}

Extinction is the regime in which the cycle ceases to turn. Value is no longer generated, circulated, or returned; the parties no longer revise their sense of one another; practice halts; and the relation decays toward the terminal rest in which it is a relation no longer. Extinction is not a benign equilibrium, a relation at peaceful rest, but the cessation of the activity by which the relation existed (\xref{p11:sec:cyc-persistence}). It is the gravest regime in the order of existence, since it is the end of the relation as such, though it is not the gravest in the order of justice, since a catastrophic cycle may inflict, in its persistence, a depletion that an extinction would have ended. The relation between extinction and the catastrophic cycle is, in this respect, genuinely tragic: the persistence that distinguishes the catastrophic cycle from extinction is also the persistence of its extraction, and the cessation that would end the extraction is also the end of the relation.

\subsection{Bifurcation between the regimes}
\label{p11:sec:reg-bifurcation}

The passage of a relation between these regimes need not be gradual. As the conditions of a relation drift, its slowly varying parameters in the sense of \xref{p11:sec:dyn-bifurcation}, its behaviour may change smoothly through a range and then, at a critical value, change in kind: a flourishing spiral may, as its recyclability declines past a threshold, lose the persistence of its ascent and tip into a catastrophic cycle or toward extinction; a catastrophic cycle may, as its depletion exhausts the conditions of its persistence, tip into extinction; and a relation near a threshold may be carried across it by a small further change in its sustaining conditions. This is the dynamical content of an observation common to the experience of relations, that their decline is often not gradual but punctuated, a long apparent stability followed by a rapid passage into a different regime. The language of bifurcation renders this: the difference between regimes is a difference in kind, and a relation may be carried from one to another by a change in degree that crosses a critical value, so that a small neglect, compounded, may tip a relation across a threshold from which its return is not symmetric to its approach.

\subsection{Justice as the internal condition of the good cycle}
\label{p11:sec:reg-justice}

The section closes by establishing the thesis toward which the regimes have been arranged: that justice is not an external constraint upon the flourishing spiral but the internal condition of its persistence with positive phase.

\begin{claim}[Justice is the internal condition of the flourishing spiral]
The flourishing spiral, the persisting cycle of positive phase, is necessarily just, and the catastrophic cycle is the form that injustice takes in a persisting relation. A cycle whose return distribution is severely skewed cannot accumulate a positive phase for all its parties, since the party consigned to the low-return tail is depleted, not developed, by the turning; its phase, for that party, is negative, and the cycle is for that party catastrophic however it appears in the aggregate. A positive phase for all the parties therefore requires a return distribution that is not severely skewed, which is to say it requires the condition of justice. Justice is, accordingly, not a constraint added to the flourishing of a relation from outside it, but the internal condition under which a cycle can ascend for all its parties rather than for some at the cost of others. This is the dynamical form of the prior paper's thesis that justice constitutes rather than constrains eudaimonia, and the skew of the return distribution (\xref{p11:sec:ev-skew}) is its statistical signature: a relation is a flourishing spiral only if its return is, relative to its coupling, approximately symmetric, and a severely skewed return is the signature of a catastrophic cycle wearing, in the aggregate, the appearance of a spiral.
\end{claim}

\noindent
The three regimes, so arranged, gather the results of the paper into a single dynamical picture. A relation lives by the persistence of its cycle and flourishes by the positive phase of its turning, and it flourishes for all its parties only if its turning is just. The instruments of evaluation proxy these features, the recyclability proxying persistence, the skew proxying the justice on which the positive phase depends, and the satisfaction measures proxying an aggregate that conceals the distinction the skew reveals; and the moments of the cycle cultivate them, the practice forming the habits that sustain the turning, the reflection guarding against the forged appearance of a spiral, and the re-cognition carrying the phase of each turn into the next. The paper now concludes, in the plural conclusions its method requires.
